% ===================================================================
%  तालिका सं. 15 — बाज़ार। — खाद्य सामग्री।
%  Traduction 2026 d'apres texto/io, controlee sur texto/fr et
%  texto/en. Consigne : voir texto/hi/10-talika-01.tex.
% ===================================================================

%%K t15-apar-1 apar
\VUsaut{4.0mm}
\VUtitre{15.0pt}{60}{तालिका सं. 15}
\VUsaut{3.0mm}

%%K t15-tit-1 sub
\VUpk{11.4pt}{40}{बाज़ार। --- खाद्य सामग्री।}
\VUsaut{3.0mm}

%%K t15-01-1 p
1. --- जो व्यापारी हमें हमारी ज़रूरत का खाना देते हैं, उनमें से अधिकतर
\VUgras{ढके बाज़ार} के \VUgras{हॉल} \textsuperscript{(1)} के नीचे जमा रहते हैं, या
पास की गलियों में।

%%K t15-02-1 p
2. --- \VUgras{सूअर के माँस का क़साई} \textsuperscript{(2)}, जो \VUgras{हैम}
\textsuperscript{(3)}, \VUgras{ख़ून की सॉसेज} \textsuperscript{(4)},
\VUgras{सॉसेज} \textsuperscript{(5)}, \VUgras{आँतों की सॉसेज} \textsuperscript{(6)}
और \VUgras{चरबी} \textsuperscript{(7)} बेचता है, वह
\VUgras{मक्खन और पनीर बेचने वाली} \textsuperscript{(8)} के पास बैठता है, जिसका
ठेला लाल छिलके वाले \VUgras{हॉलैंड के पनीर} \textsuperscript{(9)},
\VUgras{कामांबेर} \textsuperscript{(10)}, \VUgras{ग्रुयेर के बड़े चक्के}
\textsuperscript{(11)} और ताज़े \VUgras{मक्खन के पेड़ों} \textsuperscript{(12)} से
सजा है।

%%K t15-03-1 p
3. --- उसके सामने एक \VUgras{सब्ज़ी बेचने वाली} \textsuperscript{(13)} अपनी
ग्राहकों को सुंदर \VUgras{टमाटर} \textsuperscript{(14)}, \VUgras{फूलगोभी}
\textsuperscript{(15)}, \VUgras{सलाद} \textsuperscript{(16)}, \VUgras{बंदगोभी}
\textsuperscript{(17)}, \VUgras{कद्दू} \textsuperscript{(19)}, \VUgras{ख़रबूज़े}
\textsuperscript{(18)} और \VUgras{कुकुरमुत्ते} \textsuperscript{(20)} तक दिखाती है।
पर एक \VUgras{शराब वाले} ने अपनी \VUgras{ढुलाई की गाड़ी} \textsuperscript{(21)},
जो \VUgras{बीयर के पीपों} \textsuperscript{(22)} से लदी है, ठीक उसके ठेले के सामने
खड़ी कर दी है और उसकी ग्राहकों को पास नहीं आने देती। एक बाग़वान उसे
\VUgras{हाथीचक} \textsuperscript{(23)} की एक टोकरी ला रहा है।

%%K t15-tit-2 sub
\VUsaut{4.0mm}
\VUpk{10.2pt}{40}{बेचना और ख़रीदना।}
\VUsaut{3.0mm}

%%K t15-04-1 p
4. --- बाज़ार में बड़ी चहल-पहल है, क्योंकि \VUgras{नीलामी} \textsuperscript{(24)}
का समय आ गया है। जो \VUgras{गृहिणियाँ} \textsuperscript{(26)} यहाँ सौदा करने आती
हैं, वे रास्ते में \VUgras{ओझड़ी बेचने वाली} \textsuperscript{(25)} के ठेले के
सामने रुककर \VUgras{ओझड़ी} या एक \VUgras{बछड़े का सिर} ख़रीदती हैं।

%%K t15-05-1 p
5. --- एक मोटा \VUgras{रसोइया} \textsuperscript{(27)} \VUgras{मछली बेचने वाली}
\textsuperscript{(28)} से एक बहुत सुंदर \VUgras{टूना} का मोल-भाव कर रहा है, जो
अपने मन से दाम बढ़ाती जाती है। उसके यहाँ \VUgras{बाम, सोल, ट्राउट} और
\VUgras{कार्प} की बड़ी खेप आई है। उसका \VUgras{कॉड} और उसकी \VUgras{सीपियाँ}
बहुत ताज़ी हैं।

%%K t15-tit-3 sub
\VUsaut{4.0mm}
\VUpk{10.2pt}{40}{सस्ता माल और महँगा।}
\VUsaut{3.0mm}

%%K t15-06-1 p
6. --- उसके पास बैठी \VUgras{मुर्ग़ी बेचने वाली} \textsuperscript{(29)} बहुत
ईमानदार है। जब वह कोई \VUgras{टर्की}, कोई \VUgras{हंस}, कोई \VUgras{बत्तख} या
साधारण \VUgras{मुर्ग़ी} बेचती है, तो उचित दाम से अधिक नहीं माँगती।

%%K t15-07-1 p
7. --- चौक के कोने पर, पहली मंज़िल पर, एक \VUgras{कपड़े की दुकान वाली}
\textsuperscript{(30)} ने हाल ही में दुकान खोली है, जहाँ ख़रीदने वालों को
\VUgras{मेज़पोश}, \VUgras{नैपकिन, एप्रन} और \VUgras{तौलियों} का बहुत बढ़िया
चुनाव सदा मिलता है। उसी मंज़िल पर एक \VUgras{छुरी बनाने वाले}
\textsuperscript{(31)} ने अपनी कार्यशाला लगाई है। वह बाज़ार की औरतों की
\VUgras{छुरियाँ} तेज़ करता है और माली के लिए सबसे कड़े \VUgras{इस्पात} से
\VUgras{छाँटने के हँसुए} बनाता है।

%%K t15-tit-4 sub
\VUsaut{4.0mm}
\VUpk{10.2pt}{40}{परचून की दुकान।}
\VUsaut{3.0mm}

%%K t15-08-1 p
8. --- उसकी कार्यशाला के नीचे कुछ समय पहले एक बड़ी खाद्य-दुकान खुली है। यह अब
पुराने दिनों की वह सादी छोटी \VUgras{परचून की दुकान} \textsuperscript{(32)} नहीं
रही, जो केवल रोज़ का माल बेचती थी — \VUgras{चाय} \textsuperscript{(33)},
\VUgras{चॉकलेट} \textsuperscript{(34)}, \VUgras{डली की चीनी} \textsuperscript{(37)}
और \VUgras{साबुन} \textsuperscript{(38)}। यहाँ सब कुछ बिकता है :
\VUgras{कुरकुरे बिस्किट}, \VUgras{डिब्बाबंद चीज़ें} \textsuperscript{(35)},
\VUgras{मदिराएँ} \textsuperscript{(36)}, \VUgras{रम, ब्रांडी, किर्श, अनीसेत} और
\VUgras{कुरासाओ}। यहाँ शिकार भी उतनी ही सहजता से मिलता है —
\VUgras{ख़रगोश} \textsuperscript{(39)}, \VUgras{चिड़ियाँ} \textsuperscript{(40)},
\VUgras{तीतर} \textsuperscript{(41)}, \VUgras{बटेर} \textsuperscript{(42)},
\VUgras{जंगली बत्तख} \textsuperscript{(43)} या \VUgras{तीतर-मोर}
\textsuperscript{(44)} — जितनी उष्ण देशों के फल, जैसे \VUgras{केले}
\textsuperscript{(46)} और \VUgras{अनन्नास}।

%%K t15-09-1 p
9. --- एक बहुत उपकारी परचून का \VUgras{कारिंदा} \textsuperscript{(45)} जो माँगा
जाए वही देने को तैयार रहता है : \VUgras{मुरब्बा} \textsuperscript{(47)}, करौंदे का
या बिही का, \VUgras{चरबी} \textsuperscript{(48)}, \VUgras{खीरे}
\textsuperscript{(49)} या नमकीन \VUgras{सार्डिन} \textsuperscript{(50)}।

%%K t15-09-2 p
यह \VUgras{ग्राहिकाओं} \textsuperscript{(51)} के लिए बहुत सुविधा की बात है, जो
सबसे साधारण माल के पास ही सबसे दुर्लभ अगेती उपज और सबसे स्वादिष्ट फल पाती हैं :
रसीली \VUgras{नाशपाती} \textsuperscript{(52)}, \VUgras{आलूबुख़ारे}
\textsuperscript{(53)}, \VUgras{अंगूर} \textsuperscript{(54)}, \VUgras{आड़ू}
\textsuperscript{(55)}, \VUgras{बादाम} \textsuperscript{(56)} और
\VUgras{अख़रोट} \textsuperscript{(57)}।

%%K t15-tit-5 sub
\VUsaut{4.0mm}
\VUpk{10.2pt}{40}{ग्राहक।}
\VUsaut{3.0mm}

%%K t15-10-1 p
10. --- \VUgras{चावल} \textsuperscript{(58)} और \VUgras{मसूर} \textsuperscript{(59)}
धुँआई \VUgras{हेरिंग} \textsuperscript{(60)} के डिब्बे के पास रखे हैं;
\VUgras{आलू} \textsuperscript{(61)} और \VUgras{सेम} \textsuperscript{(63)}
\VUgras{सेबों} \textsuperscript{(62)}, \VUgras{अंजीरों} \textsuperscript{(64)},
\VUgras{सूखे आलूबुख़ारों} \textsuperscript{(65)} और \VUgras{पहाड़ी बादामों}
\textsuperscript{(66)} के पास बैठे हैं। इस दुकान में ताज़े \VUgras{अंडे}
\textsuperscript{(67)} और \VUgras{जैतून} \textsuperscript{(68)} भी मिलते हैं, सो
\VUgras{टोकरियाँ} \textsuperscript{(70)} और \VUgras{जाली के थैले}
\textsuperscript{(73)} बहुत जल्दी भर जाते हैं।

%%K t15-11-1 p
11. --- इस बढ़िया दुकान की \VUgras{कॉफ़ी} \textsuperscript{(72)} ने
\VUgras{नौकरानियों} \textsuperscript{(69)} और रसोइयों में अच्छा-ख़ासा नाम कमाया
है, क्योंकि बूढ़ा \VUgras{परचून वाला} \textsuperscript{(71)} उसे स्वयं बड़े ध्यान
से भूनता है।

%%K t15-tit-6 sub
\VUsaut{4.0mm}
\VUpk{10.2pt}{40}{देखने की चीज़ें।}
\VUsaut{3.0mm}

%%K t15-12-1 p
12. --- सब लोग बाज़ार सौदा करने ही नहीं आते। हॉल तक जाती उस छोटी गली की सुरम्य
आवाजाही में तरह-तरह के लोग दिखाई देते हैं। एक भले \VUgras{कसाईख़ाने के मज़दूर}
\textsuperscript{(75)} के पास, जो अपने आगे एक \VUgras{ठेला} \textsuperscript{(74)}
धकेल रहा है, यह रहे दो छुट्टी पर आए सिपाही, जिन्हें अपनी चमकीली वरदी दिखाने में
बड़ा गर्व है : \VUgras{हुसार} \textsuperscript{(76)} अपने नीले \VUgras{दोलमान} और
\VUgras{पंखदार शाको} में, और \VUgras{ज़ुआव का हवलदार} ढीली निकर में, सिर पर
\VUgras{फ़ेज़} लगाए।

%%K t15-13-1 p
13. --- \VUgras{नानबाई} \textsuperscript{(80)}, जो \VUgras{नानबाईख़ाने}
\textsuperscript{(78)} की देहरी पर खड़ा है, उसने अभी-अभी अपनी \VUgras{रोटी}
\textsuperscript{(79)} का आटा गूँधा है और तंदूर से वे \VUgras{क्रोइसाँ}
\textsuperscript{(81)}, \VUgras{छोटी रोटियाँ} \textsuperscript{(82)} और
\VUgras{गोल रोटियाँ} \textsuperscript{(83)} निकाली हैं जो खिड़की में रखी हैं।

%%K t15-14-1 p
14. --- नानबाईख़ाने के ऊपर एक कुशल \VUgras{सिलाई करने वाली} \textsuperscript{(84)}
रहती है, जो कई \VUgras{कढ़ाई करने वालियाँ} \textsuperscript{(85)} रखती है। उसके
पास एक \VUgras{धोबिन और इस्तिरी करने वाली} \textsuperscript{(86)} रहती है, जो
\VUgras{कपड़े} \textsuperscript{(88)} धोकर और सुखाकर उनमें कलफ़ लगाती है और
\VUgras{इस्तिरी} \textsuperscript{(87)} से दबाती है।

%%K t15-tit-7 sub
\VUsaut{4.0mm}
\VUpk{10.2pt}{40}{माँस, मछली और सब्ज़ी।}
\VUsaut{3.0mm}

%%K t15-15-1 p
15. --- भूतल की \VUgras{क़साई की दुकान} \textsuperscript{(89)} में हर तरह का माँस
बिकता है — \VUgras{भेड़ का} \textsuperscript{(90)}, \VUgras{गाय का}
\textsuperscript{(94)} या \VUgras{बछड़े का} \textsuperscript{(92)} — \VUgras{रानें,
चापें, टुकड़े} और \VUgras{पसलियाँ} के रूप में। \VUgras{क़साई का छोकरा}
\textsuperscript{(93)} यह सारा माँस काटता है और \VUgras{मज्जा की हड्डियाँ}
\textsuperscript{(96)} अलग रखता है। क़साई के दरवाज़े पर एक
\VUgras{सीपी बेचने वाली} \textsuperscript{(97)} खड़ी है, जो \VUgras{सीपियाँ}
\textsuperscript{(98)} बेचती है। चाहें तो वह उन्हें खोल भी देती है।

%%K t15-16-1 p
16. --- ये सब व्यापारी लाइसेंस या जगह का शुल्क देते हैं; इसीलिए
\VUgras{सिपाही} \textsuperscript{(100)} उन बेचारी \VUgras{घूमकर सब्ज़ी बेचने वालियों}
\textsuperscript{(99)} का बिना दया पीछा करता है, जो अपनी ठेलियों पर
\VUgras{गाजर, मूली, शलजम, गंधना} या \VUgras{शतावरी के गट्ठर, ताज़े मटर, पालक,
चूका, जलकुंभी} और \VUgras{धनिया} ढोकर उनसे प्रतिस्पर्धा करती हैं।

%%K t15-tit-8 sub
\VUsaut{4.0mm}
\VUpk{10.2pt}{40}{सिपाही से बचो!}
\VUsaut{3.0mm}

%%K t15-17-1 p
17. --- \VUgras{लहसुन} \textsuperscript{(101)} और \VUgras{प्याज़} बेचने वाला लड़का
ख़ूब जानता है कि उसे भी बाज़ार के पास अपना माल बेचने की अनुमति नहीं। वह भाग खड़ा
होता है, इस ख़तरे के साथ कि \VUgras{झींगे-मछली} और \VUgras{लॉब्स्टर}
\VUgras{बेचने वाली} \textsuperscript{(102)} की \VUgras{केकड़ों},
\VUgras{क्रेफ़िश} और \VUgras{झींगों} की टोकरी उलट दे।

%%K t15-18-1 p
18. --- सब लोग हड़बड़ी में हैं और बड़ा शोर मचा रहे हैं; पर इससे घूमते
\VUgras{शीशासाज़} \textsuperscript{(103)} की हाँक साफ़ सुनाई देना नहीं रुकता, न उस
\VUgras{कोयले वाले} \textsuperscript{(104)} की, जिसने अभी-अभी एक
\VUgras{कोयले की बोरी} \textsuperscript{(105)} पहुँचाने को अपनी गाड़ी थोड़ी देर के
लिए छोड़ी है।
\VUsaut{6.0mm}
\VUfilet{22mm}
